% BHU PYQ -- Professional Exam Template
\documentclass[11pt,a4paper]{article}

\usepackage[a4paper,top=2.5cm,bottom=2.5cm,left=2.8cm,right=2.8cm,headheight=14pt]{geometry}
\usepackage{amsmath,amssymb}
\usepackage[T1]{fontenc}
\usepackage[utf8]{inputenc}
\usepackage{lmodern}
\usepackage{microtype}
\usepackage{fancyhdr}
\usepackage{enumitem}
\usepackage{listings}
\usepackage{hyperref}
\usepackage{lastpage}
\usepackage{parskip}

\hypersetup{colorlinks=true,urlcolor=black,linkcolor=black,
  pdftitle={CS-109: Data Communication -- BHU PYQ}}

\pagestyle{fancy}
\fancyhf{}
\renewcommand{\headrulewidth}{0.4pt}
\renewcommand{\footrulewidth}{0.4pt}
\fancyhead[L]{\small   \textit{Banaras Hindu University}}
\fancyhead[R]{\small   \textit{CS-109: Data Communication}}
\fancyfoot[C]{\small Page  \thepage\ of \pageref{LastPage}}

\lstset{
  basicstyle=\small tfamily,
  frame=single,
  breaklines=true,
  columns=flexible,
  keepspaces=true
}


\newlist{parts}{enumerate}{2}
\setlist[parts,1]{label=(\alph*),leftmargin=*,itemsep=4pt,topsep=3pt}
\setlist[parts,2]{label=(\roman*),leftmargin=*,itemsep=2pt,topsep=2pt}

\newcommand{\pts}[1]{\hfill   \textit{[#1]}}


\begin{document}

\begin{center}
  {\large\bfseries BANARAS HINDU UNIVERSITY}\\[2pt]
  {
\normalsize B.Sc. (Hons.) Semester V Examination, 2022-23}\\[6pt]
  \rule{0.8\linewidth}{0.5pt}\\[6pt]
  {\large\bfseries Paper: CS-109 --- Data Communication}\\[4pt]
  {
\normalsize Time: Three hours \hfill Maximum Marks: 70}
\end{center}

\rule{\linewidth}{0.4pt}

\smallskip



\noindent   \textit{Instructions: Attempt any five questions. Figures on the right-hand side margin indicate max}

\bigskip

\medskip


\noindent   \textbf{Question 1.} (a) Define a System. What are the key elements of a system? Distinguish between Conceptual, Physical and Logical System Models.\pts{7}

\begin{parts}
  \item Explain through a neat diagram the hierarchy of information levels and information systems in an organization.\pts{7}
\end{parts}

\medskip


\noindent   \textbf{Question 2.} (a) Explain the Software Development Life Cycle (SDLC) model along with its major advantages and disadvantages.\pts{7}

\begin{parts}
  \item Explain the commonly used Fact Finding Tools for Preliminary Investigation.\pts{7}
\end{parts}

\medskip


\noindent   \textbf{Question 3.} (a) What is Feasibility Study? Explain the various types of feasibility analysis carried out\pts{7}
during system analysis.

\begin{parts}
  \item What is the purpose of designing SRS document? Explain the typical structuring\pts{7}
  elements of SRS document.
\end{parts}

\medskip


\noindent   \textbf{Question 4.} (a) Distinguish between Coupling and Cohesion as design principles. What is desirable\pts{7}
and why?

\begin{parts}
  \item Describe the following terms in reference to System testing:\pts{7}
  Error, Fault, Failure, Reliability
\end{parts}

\medskip


\noindent   \textbf{Question 5.} (a) Distinguish between verification \& Validation. What are the various levels of\pts{7}
software system testing?

\begin{parts}
  \item What is the major difference between Blackbox and Whitebox testing? Explain the\pts{7}
  Various criteria for testcase generation in the two approaches.
\end{parts}

\medskip


\noindent   \textbf{Question 6.} (a) Describe the commonly used methods for System Conversion.\pts{7}

\begin{parts}
  \item Differentiate between Corrective, Adaptive and Perfective maintenance of a system.\pts{7}
\end{parts}

\medskip


\noindent   \textbf{Question 7.} Write short notes on any two of the following:\pts{14}

\begin{parts}
  \item Tangible and Intangible Costs and Benefits
  \item Structured Analysis
  \item Role of a System Analyst
  \item Software Cost Estimation
\end{parts}

\vfill
\rule{\linewidth}{0.4pt}

\begin{center}\small
\href{https://bhu-exam.vercel.app/}{Click here to visit our website}\\[4pt]\href{https://me-aryan.vercel.app/}{Click here to visit our contributor}\\[4pt]\href{https://quashan-me.vercel.app/}{Click here to visit our contributor}
\end{center}
\end{document}